%% notes-tex/database/database-expansion-hosts/sections/state.tex
%% [[experiments.database.expansion-hosts]]
%% https://github.com/Mjvolk3/torchcell/tree/main/notes-tex/database/database-expansion-hosts/sections/state.tex
%% SOURCE: hand-authored prose. Served dataset and record counts come from release
%% 2026.10.17-4295f7ca and the per-organism tables of the thesis database chapter
%% (experiments/database/results/organism_summary.json, probe variant probes-v1); the
%% candidate figures come from the triage report
%% /scratch/projects/torchcell-scratch/tmp/deep-research/reports/Pk and next dataset wave triage.md
%% and the research notes beside it.

\section{State across hosts\secstatus{todo}}\label{sec:state}

A \term{host} here is an organism the served graph can declare, meaning one that
\texttt{torchcell/data/experiment\_dataset.py} holds as an \texttt{Organism} member and the
genomes tier resolves to pinned bytes. Three hosts are served today and the asymmetry between
them is the starting condition of this pass: \org{S. cerevisiae} carries 70 datasets and
65,889,109 records, \org{E. coli} 68 datasets and 4,561,865 records, and \org{P. putida} 32
datasets and 2,695,603 records. The three-organism total of 73,146,575 is computed against a
yeast figure of 65,889,107, two records below the figure the yeast table itself reports, so
the discrepancy is in the served tables and is carried here rather than resolved.

Where each host is thin is a property of its phenotype mix rather than of its record count.
\org{E. coli}'s largest section is environmental and chemogenomic response at 22 datasets,
followed by fitness, metabolite, and protein abundance at 8 each, expression and viability at
5 each, product titer at 3, and flux, gene interaction, and protein synthesis and turnover at
2 each, with one dataset each for morphology, phosphoproteome, and promoter activity. For
\org{P. putida} protein abundance leads with 14 datasets, then environmental and chemogenomic
response with 7, product titer with 6, metabolite with 2, and expression, fitness, and
viability with one each. So the newer host of the two carries almost no gene-perturbation
fitness, which is the layer the yeast corpus is built on, and the layer a transfer claim would
be judged against.

The plan this serves is stated in the thesis conclusions chapter: ``The first pillar is the
database, which already holds \org{Escherichia coli} and extends next to ten further
gram-negative bacteria and ten yeasts, each with its own virtual cell, with \org{Pseudomonas
putida} and \org{Pichia kudriavzevii} as the first hosts with little data, on which the models
are judged.'' The same passage records the test as a comparison of label counts. A Cell Graph
Transformer trained on
\org{S. cerevisiae} and fine-tuned on \org{P. kudriavzevii}, both budding yeasts, or trained
on \org{E. coli} and fine-tuned on \org{P. putida}, both in the class Gammaproteobacteria,
could match a model trained from scratch on the target with several-fold fewer labels, and the
threshold recorded beside it is five-fold fewer labels at matched Pearson correlation.

That test fixes what a low-data host needs before a model can be judged on it, and the list is
short and ordered. It needs a reference genome pinned in the genomes tier, and
\org{P. kudriavzevii} SD108 already has one, because
\file{torchcell/sequence/genome/registry.py} pins the assembly set
\texttt{pkudriavzevii\_SD108\_v1} and \file{torchcell/candidates/gates.py} already lists that
set among the readable ones. It needs an \texttt{Organism} member, which that host
does not have, so no loader can declare it and \texttt{dataset\_adapter\_map} can hold nothing
for it. It needs a host token in the gate, so that G2 can resolve a row's strain to that
assembly set. And it needs held-out labels of the same kind as the source host's labels, and
the subsections below measure how far each host of this wave is from that.

\subsection{What the served graph holds per host\secstatus{todo}}\label{sec:state-served}

On the gene-perturbation fitness axis the served graph is one host deep. The yeast corpus's
fitness and gene-interaction maps supply 67 percent of all 65,889,109 yeast records, and the
metabolite and titer screens that bear directly on strain engineering are three to four orders
of magnitude smaller. \org{P. putida} contributes one fitness dataset and \org{E. coli} eight,
so a cross-host fitness comparison has a dense source and two sparse targets.

On the sequence-to-readout axis the served graph is thinner still, on every host. One dataset
serves a varied-sequence library at scale, \texttt{vaishnav2022.py} with 71,158 records of
\texttt{VariablePromoterReporterPerturbation} crossed with
\texttt{PooledSortSeqExpressionPhenotype}, and \org{E. coli} carries a single promoter-activity
dataset. Everything else perturbs a gene list. That matters for a low-data host because a
promoter, a terminator, and an integration locus are the parts a new host is engineered with
first, before any deletion collection for it exists.

\subsection{What this wave adds per host\secstatus{todo}}\label{sec:state-rows}

\notesec{\org{S. cerevisiae}, measurements rather than rows} The three regulatory-DNA rows add
102,371,025, 51,339,035 and 6,739,258 measured promoter sequences, all on the same sort-seq
platform and all under classes the served schema already holds (Sec.~\ref{sec:regdna}). Below
them the same axis supplies a 1,000,000-sequence transcription-start-site library, a
489,348-sequence 5' UTR library, a 13,000-sequence 3' end library, and the position-effect
panels, each of which needs one new leaf rather than a new record type.

\notesec{\org{P. putida}, two ready-made matrices} The additions are dense and small in row
count: 332 RB-TnSeq samples as a 4,732 by 332 gene fitness matrix with t statistics and
per-sample metadata, 321 RNA-seq profiles as a 5,564 by 321 matrix over 118 typed conditions,
a 71-condition nitrogen BarSeq screen with 672 genes showing a significant phenotype, and a
24-strain by 205-substrate respiration panel that is the only resource letting this host enter
as a species with intraspecific variation rather than as one reference strain
(Sec.~\ref{sec:pseudo}).

\notesec{Nine further pseudomonads, inside one artifact} The nine non-KT2440 pseudomonad rows
of the Fitness Browser hold 1,479 genome-wide fitness experiments between them, measured, and
roughly $6.7\times10^{6}$ measurements, inferred from the per-gene matrix shape. They arrive
in the same \texttt{feba.db} download, under the same schema, with the same condition
vocabulary and the same RefSeq mapping files as the KT2440 rows already served, so the
marginal work is genome registration rather than loader work.

\notesec{\org{P. aeruginosa}, essentiality and isolate panels} This host enters for the
layer the others lack: 199 curated essentiality calls scoped by strain and medium, a
414-isolate clinical panel with a genome, a transcriptome, and 1,656 agar-dilution MIC records
per the CLSI standard, a 443-isolate longitudinal panel with 8 phenotypes, and a
125-isolate survey with 13 readouts (Sec.~\ref{sec:pseudo}).

\notesec{\org{P. kudriavzevii}, about 25 small rows} The host's additions are a 161-isolate by
57-stressor growth matrix, 9,177 measurements inferred, a 12-strain low-pH expression atlas at
two sampling points, 128 four-drug MIC values on 32 sequenced isolates, a 47-guide by
21-locus landing-pad screen with four readouts, and four small GEO series
(Sec.~\ref{sec:pk}). None of them is a screen.

\subsection{What is still missing after this wave\secstatus{todo}}\label{sec:state-gaps}

Four gaps survive the wave, and each is a measured absence rather than an unexamined axis.

There is no measured essentiality or genome-wide fitness screen in \org{P. kudriavzevii} of
any kind, so a model trained on yeast essentiality has no target to be judged against in that
host. Its only measured gene-perturbation fitness data are the 47-guide landing-pad screen,
one validated paralog-pair deletion, and the targeted knockouts inside each engineering paper.
Essentiality exists only as output of the iIsor850 reconstruction over 850 genes, which the
yeast pool already excludes under the rule \texttt{not-a-dataset}. The one remaining lead is
unresolved rather than negative: NCBI holds 713 BioProjects for the species and 357 of them
match \texttt{resequencing}, and three sampled ones registered by the DOE Joint Genome
Institute on 2026-05-11 are titled as SD108 \texttt{IoSL} resequencing with no publication
attached.

There is no baseline genome-wide essentiality call for KT2440 from the 2018 to 2026 window at
all, because RB-TnSeq by construction cannot score a gene it cannot mutate. The dataset that
would complete it, a 78,932-member CRISPRi library with 10 to 15 guides per gene and 798
non-targeting controls screened on glucose, acetate, and p-coumarate at two concentrations,
exists only as a DOE abstract with no publication, no repository, and no accession, so it is
tracked rather than planned for.

There is no \org{P. putida} or environmental-pseudomonad isolate collection with paired
sequence and growth data beyond the 24-strain respiration panel, and that panel's
data-availability statement could not be read. So the intraspecific genotype axis for this host
rests on one preprint.

And there is no ribosome profiling, CAGE, transcription-start-site mapping, or arrayed
overexpression dataset for \org{P. kudriavzevii}, by a targeted query and the GEO taxid sweep.
Chinese-language indices, ArrayExpress and BioStudies, and the Zenodo, Dryad, and figshare
APIs were not searched, so absence on those routes is unmeasured rather than negative.
